\section{Literature Review and Research Gap}

\subsection{From Trust to Architecture}

The foundational trust-in-automation literature established that performance in sociotechnical systems depends on appropriate reliance rather than simple acceptance or rejection \citep{lee2004trust}. This insight remains indispensable, but its typical unit of analysis is a particular automation interaction or local reliance decision. In AI-enabled organizations, however, actors also decide where AI should sit in the workflow, how much authority it should receive, and what forms of override or confirmation should remain in place. These are architecture questions, not merely trust questions.

\subsection{Algorithm Aversion, Appreciation, and the Limits of One-shot Framing}

Research on algorithm aversion demonstrates that people may avoid algorithms after observing them err, even when the algorithm performs better overall \citep{dietvorst2015algorithm}. By contrast, algorithm appreciation research shows that people may prefer algorithmic judgment under some circumstances \citep{logg2019algorithm}. Together, these findings imply that humans do not evaluate AI output in a source-neutral way. Still, the dominant framing remains one-shot: whether an algorithm is accepted, weighted, or rejected at a given moment. The present project asks a stronger question: can repeated asymmetries in how AI errors are learned produce long-run convergence toward low-authorization architectures?

\subsection{Confidence Calibration Is Not Role Learning}

Another important stream studies how confidence scores, explanations, and uncertainty signals affect trust calibration in AI-assisted decision making \citep{zhang2020confidence,ma2023trust}. This literature is chiefly concerned with signal-level learning: whether users learn the mapping between case-specific AI signals and correctness likelihood. That question is analytically distinct from the one pursued here. Even perfectly calibrated confidence signals do not by themselves specify whether AI should act autonomously, advise a human, or require explicit confirmation. The present research therefore differentiates signal-level calibration from architecture-level learning.

\subsection{Human--AI Teaming and Learning to Defer}

Recent research in machine learning and HCI has begun to consider how human and AI agents should divide labor. Learning-to-defer approaches ask when models should pass a decision to human experts \citep{mozannar2020defer}, and studies of human--AI teams examine how complementarity, performance tradeoffs, and updating affect joint outcomes \citep{bansal2019updates,vaccaro2024combinations}. These contributions are highly relevant, but they often approach the allocation problem from the system or model side. By contrast, this project foregrounds the human side of architecture formation: how people infer an appropriate role for AI under repeated feedback.

\subsection{EWA and Path Dependence as Integrating Mechanisms}

EWA provides a natural dynamic language for the proposed argument because it explains how repeated experience updates the attraction of alternative strategies \citep{camerer1999ewa}. Path dependence research, meanwhile, highlights how early events may shape long-run outcomes through lock-in and increasing returns \citep{david1985qwerty,arthur1989lockin}. Combining these traditions suggests a new mechanism for human--AI collaboration: early AI errors can induce stronger negative updating, leading actors to choose lower-autonomy structures, which in turn reduce future exposure to potentially trust-repairing AI performance, thereby slowing recovery and stabilizing low-AI-authorization equilibria in the loose behavioral sense.

\subsection{The Core Gap}

The literature therefore leaves three linked questions underdeveloped. First, how do humans learn AI's role, not just AI's reliability? Second, can asymmetric reactions to AI versus human errors shape long-run authorization structures rather than immediate preferences alone? Third, what formal mechanism can connect repeated feedback, architecture attraction, path dependence, and behavioral steady states? The proposed project is designed to answer these questions.
